\documentclass[10pt,a4paper]{amsart}
\usepackage[margin=19mm]{geometry}
\usepackage{amsmath,amssymb,mathtools}
\usepackage[scr=rsfs,cal=euler]{mathalfa}
\usepackage{xfrac}
\usepackage[unicode,hidelinks,bookmarks=false]{hyperref}
\newcommand{\ie}{i.e.\ }
\newcommand{\cf}{cf.\ }
\newcommand{\resp}{respectively\ }
\newcommand{\Subsection}{\subsection}
\providecommand{\nobreakdash}{\nobreak\hbox{-}}
\setlength{\emergencystretch}{2em}
\multlinegap0pt
\usepackage{bm}
\usepackage{bbm}
\usepackage{dsfont}
\usepackage{xspace}
\usepackage{cancel}
\usepackage{mathtools}
\usepackage{cleveref}
\usepackage[shortlabels]{enumitem}
\usepackage{amsthm}
\makeatletter
\@ifundefined{theorem}{\newtheorem{theorem}{Theorem}}{}
\@ifundefined{proposition}{\newtheorem{proposition}[theorem]{Proposition}}{}
\makeatother
\DeclareMathOperator{\C}{\mathcal{C}}
\newcommand{\F}{{\mathbb F}}
\newcommand{\GL}{\hbox{{\rm GL}}}
\newcommand{\M}{\mathrm{M}}
\newcommand{\Tr}{{\mathrm{Tr}}}
\title[Adjoint and duality for rank-metric codes in a skew polynomi]{Adjoint and duality for rank-metric codes in a skew polynomial framework}
\author{Jos\'e G\'omez-Torrecillas, F. J. Lobillo, Gabriel Navarro, Paolo Santonastaso}
\date{Source: arXiv:2511.05084v1 (CC BY 4.0)}
\begin{document}
\maketitle

\noindent\emph{Source and licence.} Adjoint and duality for rank-metric codes in a skew polynomial framework. Version arXiv:2511.05084v1 (CC BY 4.0), distributed under the Creative Commons Attribution 4.0 International licence, as recorded in the arXiv licence metadata for this exact version and shown on the abstract page https://arxiv.org/abs/2511.05084v1. Full rights notice: licenses/arxiv-2511.05084-CC-BY-4.0.txt.\par\smallskip

\noindent\textit{Editorial scope.} Counted unit: the Proposition whose source label is \texttt{prop:determinantiondual} in the article (source0.tex lines 716--728, source label \texttt{prop:determinantiondual}), delivered together with the article's own complete original proof of it (source0.tex lines 730--802, one \texttt{proof} environment of 73 lines). No mathematics is written, repaired or re-derived: statement and proof are the article's own text line for line. Required context retained uncounted: eq:bilinearmatrix (lines 457--459); eq:bilinearoverRF (lines 474--476); th:correspondencebilinear (lines 488--495); eq:dimensionalformula (lines 507--509); prop:determinantionadjoint (lines 625--637); adjointD (lines 628--636); adjointS (lines 628--636); eq:relationduals (lines 711--712). Every definition, display and label of those lines is retained unchanged. Omitted from this package and not redistributed: 1--456, 460--473, 477--487, 496--506, 510--624, 637--710, 713--715, 729--729, 803--1337. Nothing inside a retained excerpt is omitted or altered: every retained body is delivered exactly as the article's lines are written, so no omission receipt of any kind accompanies this package. Citations: the retained excerpts carry no citation command at all, so no bibliography entry of the article is transcribed and no reference is redistributed. Reference adaptation: none was needed and none was made; every reference inside the delivered unit resolves inside it, so no unresolved reference or citation is produced. Printed slips kept literal: no printed word, symbol, hypothesis or number of the counted statement or of its proof was altered. Layout and class: the article's own class is \texttt{article} with packages including geometry, enumerate, amsmath, amssymb, latexsym, amsthm, color, cancel, graphicx, cite, changes, lineno, hyperref, comment; the wrapper is the lane's pinned \texttt{amsart} editorial wrapper at 10pt on A4, which restates the theorem-like environments the retained text needs and adds the article's own macro definitions that the retained text needs (\texttt{\textbackslash C}, \texttt{\textbackslash F}, \texttt{\textbackslash GL}, \texttt{\textbackslash M}, \texttt{\textbackslash Tr}), with extra wrapper packages (bm, bbm, dsfont, xspace, cancel, mathtools, cleveref, enumitem). The delivered numbering is the unit's own. Assets: no retained span contains a figure environment, a graphics-inclusion command, a PDF-page inclusion, a table or a TikZ picture; no image file is copied into this package, the package declares no asset dependency and no image file is redistributed with it. Licence: version arXiv:2511.05084v1 of this article is distributed under the Creative Commons Attribution 4.0 International licence, as recorded in the arXiv licence metadata for this exact version and shown on the abstract page https://arxiv.org/abs/2511.05084v1. A full rights notice is delivered at licenses/arxiv-2511.05084-CC-BY-4.0.txt. Characters: every retained excerpt is pure ASCII, so the delivered engine, XeLaTeX with its default Latin Modern text font, reports no missing character.
\begin{equation} \label{eq:bilinearmatrix}
\langle A,B \rangle=\Tr_{q^s/p}(\Tr(AB)),
\end{equation}

    \begin{equation} \label{eq:bilinearoverRF}
    \langle a,b \rangle_F=\Tr_{q^n/p}\left((ab)_0\right)%=\Tr_{q^n/p}\left(a_0b_0-\sum_{i=1}^{sn-1}a_{sn-i}\sigma^{sn-i}(b_i)F_0\right).
     \end{equation}

\begin{theorem} \label{th:correspondencebilinear}
Let $\langle -,- \rangle$ be the bilinear form defined as in \eqref{eq:bilinearmatrix} over $M_n(\F_{q^s})$ and let $\langle -, -\rangle_F$ defined as in \eqref{eq:bilinearoverRF} over $R_F$.
Then there exists an invertible element $U \in \GL_n(\F_{q^s})$ such that 
\[
\langle M_{R_F}(a), M_{R_F}(b) U  \rangle= \langle a,b \rangle_F,
\]
for every $a,b \in R_F$.
\end{theorem}

\begin{equation} \label{eq:dimensionalformula}
\dim_K V = \dim_K W + \dim_K W^\perp.
\end{equation}

\begin{proposition} \label{prop:determinantionadjoint}
Let $\hat{F}(y)$ be the monic reciprocal polynomial of the irreducible polynomial $F(y) \in \F_q[y]$. For any $1 \leq k<n$, the following hold. 

\begin{enumerate}[I)]
\item\label{adjointS} The adjoint code of $S_{n,s,k}(\eta,\rho,F) \subseteq R_F \cong M_n(\F_{q^s})$ is equivalent to
\[
S_{n,s,k}(\rho^{-1}(\eta^{-1}),\rho^{-1}\circ \sigma^{ks},\hat{F}) \subseteq R_{\hat{F}} \cong M_n(\F_{q^s}).
\]
\item\label{adjointD} The adjoint code of  $D_{n,s,k}(\gamma,F) \subseteq R_F \cong M_n(\F_{q^s})$ is equivalent to
 \[D_{n,s,k}\left(\sigma^{s(n-k)}\left(\frac{1}{\gamma}\right),\hat{F}\right)\subseteq R_{\hat{F}} \cong M_n(\F_{q^s}).
 \]
\end{enumerate}
\end{proposition}

\begin{equation} \label{eq:relationduals}
(\C^{\perp'})^{\top} = \C^{\perp}.\end{equation} Thus, to determine the dual of a rank-metric code, we just need to compute the adjoint of the dual code of $\C$ with respect to the Frobenius bilinear form $\langle -,- \rangle$.

\begin{proposition} \label{prop:determinantiondual}
 Let $\hat{F}(y)$ be the monic reciprocal polynomial of $F(y)$. For any $1 \leq k<n$, the following hold. 

\begin{enumerate}[I)]
    \item\label{dualS} The dual code of $S_{n,s,k}(\eta,\rho,F)$ in $R_F \cong M_n(\F_{q^s})$ is equivalent to
 \[
 S_{n,s,n-k}(\rho^{-1}(\eta F_0), \rho^{-1},\hat{F}) \subseteq R_{\hat{F}} \cong M_n(\F_{q^s}).
 \]
 \item\label{dualD} The dual code of  $D_{n,s,k}(\gamma,F)$  in $R_F \cong M_n(\F_{q^s})$ is equivalent to
 \[D_{n,s,n-k}(\sigma^{sk}(\gamma),\hat{F}) \subseteq R_{\hat{F}} \cong M_n(\F_{q^s}).
 \]
 \end{enumerate}
\end{proposition}

\begin{proof}
As in the proof of \Cref{prop:determinantionadjoint}, we fix $M_{R_F}:R_F \rightarrow M_n(\F_{q^s})$ to be an $\F_{q^s}$-algebra isomorphism. Let $\langle -, - \rangle_F$ be the bilinear form defined as in \eqref{eq:bilinearoverRF} over $R_F$. \\
\emph{\ref{dualS})}  Let $S=S_{n,s,k}(\eta,\rho,F)$, we start by computing the dual $S^{\perp}$ of $S$ with respect to the bilinear form $\langle -,- \rangle_F$ of $R_F$, i.e. 
\[
S^{\perp}=\{b \in R_F \colon \langle a,b \rangle_F=0, \mbox{ for every }a \in S\}.
\]
Clearly, every monomial $\alpha x^i+RF(x^n)$, with $i \in \{1,\ldots,s(n-k)-1\}$ is orthogonal to the elements of $S$. Moreover, for any $\alpha \in \F_{q^n}$, we have that \[c = \rho^{-1}(\eta F_0\sigma^{sk}(\alpha))+\alpha x^{s(n-k)}+RF(x^n)\] is orthogonal to any element of $S$. Indeed, if $a=\sum_{i=0}^{sk-1}a_ix^i+\eta \rho(a_0)x^{sk} + RF(x^n)$ we have 
\[
\begin{split}
\langle a, c \rangle_F &= \Tr_{q^n/p}\left( \rho^{-1}(\eta F_0\sigma^{sk}(\alpha))a_0-\sigma^{sk}(\alpha) \eta \rho(a_0)F_0  \right) \\
&= \Tr_{q^n/p}\left( \rho^{-1}(\eta F_0\sigma^{sk}(\alpha))a_0)-\Tr_{q^n/p}(\sigma^{sk}(\alpha) \eta \rho(a_0)F_0  \right) \\ 
&= \Tr_{q^n/p}\left(\rho( \rho^{-1}(\eta F_0\sigma^{sk}(\alpha))a_0))-\Tr_{q^n/p}(\sigma^{sk}(\alpha) \eta \rho(a_0)F_0  \right)  \\
&= \Tr_{q^n/p}\left( \eta F_0\sigma^{sk}(\alpha)\rho(a_0)\right) - \Tr_{q^n/p}\left(\sigma^{sk}(\alpha)  \eta \rho(a_0)  F_0 \right) \\
&= 0
\end{split}
\]
As a consequence, \[
\begin{split}
S' &= \left\{\rho^{-1}(\eta F_0)\rho^{-1}(\sigma^{sk}(\alpha))+\sum\limits_{i=1}^{s(n-k)-1}a_ix^i+\alpha x^{s(n-k)}+RF(x^n), a_i \in \F_{q^s}\right\}\\
&= \left\{b+\sum\limits_{i=1}^{s(n-k)-1}a_ix^i+\sigma^{-sk}(\eta^{-1}F_0^{-1}) \sigma^{-sk}(\rho(b)) x^{s(n-k)}+RF(x^n), a_i \in \F_{q^s}\right\} \\
&= S_{n,s,n-k}(\sigma^{-sk}(\eta^{-1}F_0^{-1}), \sigma^{-sk} \circ \rho,F)
\end{split}
\]
is contained in $S^{\perp}$. Since $\langle -,- \rangle_F$ is a bilinear non degenerate  form, by \eqref{eq:dimensionalformula}, we have  \[
\dim_{\F_p}(S^{\perp})=n^2se-\dim_{\F_p}(S^{\perp})
\]
So, we get
\[
\lvert S' \rvert= \lvert S_{n,s,n-k}(\sigma^{-sk}(\eta^{-1}F_0^{-1}), \sigma^{-sk} \circ \rho,F) \rvert=q^{ns(n-k)}=\lvert S^{\perp} \rvert
\]
Thus, $S'=S^{\perp}$.

Now, let 
\(
\C=\{M_{R_F}(a) \colon a \in S\} \subseteq M_n(\F_{q^s}).
\)
We start by determine $\mathcal{C}^{\perp'}$, i.e. the dual of $\C$ with respect of $\langle -, - \rangle$. Then, by computing the adjoint of $\mathcal{C}^{\perp'}$ and by the relation in \eqref{eq:relationduals}, we will obtain $\C^{\perp}$. We know, by \Cref{th:correspondencebilinear}, there exists an invertible element $U \in \GL_n(\F_{q^s})$ such that 
\[
\langle M_{R_F}(a), M_{R_F}(b) U  \rangle= \langle a,b \rangle_F,
\]
for all \(a,b \in R_F\). As a consequence, we have
\[
\begin{split}
\C^{\perp'} &= \{\M_{R_F}(b)U: \langle a,b   \rangle_F=0, \mbox{ for every }a \in R_F\} \\
&= \{\M_{R_F}(b): \langle a,b  \rangle_F=0, \mbox{ for every } a \in R_F\}U \\
&= \{\M_{R_F}(b): b \in S_{n,s,n-k}(\sigma^{-sk}(\eta^{-1}F_0), \sigma^{-sk} \circ \rho,F) \}U
\end{split}
\]

Therefore, $\C^{\perp'}$ is equivalent to $S_{n,s,n-k}(\sigma^{-sk}(\eta^{-1}F_0^{-1}), \sigma^{-sk} \circ \rho,F)$ in $R_F \cong M_n(\F_{q^s})$. Finally, $\C^{\perp}=(\C^{\perp'})^{\top}$ is equivalent to to the adjoint of $S_{n,s,n-k}(\sigma^{-sk}(\eta^{-1}F_0), \sigma^{-sk} \circ \rho,F)$ that is
\[
S_{n,s,n-k}(\rho^{-1}(\eta F_0),  \rho^{-1},\hat{F}) \subseteq R_{\hat{F}} \cong M_n(\F_{q^s})
\]
by \emph{\ref{adjointS})} of \Cref{prop:determinantionadjoint}, that proves our assertion. \\
\emph{\ref{dualD})} We argue as in \emph{\ref{dualS})}. Let $D=D_{n,s,k}(\gamma,F)$. We start by computing the dual $D^{\perp}$ of $D$ with respect to the bilinear form $\langle -,- \rangle_F$ of $R_F$. 
Clearly, every monomial $\alpha x^i+RF(x^n)$, with $i \in \{1,\ldots,s(n-k)-1\}$ are orthogonal to the elements of $D$. Let now $\zeta \in \F_{q^n}^*$ be a nonzero element such that $\Tr_{q^n/q^{n/2}}(\zeta \gamma)=0$. It is easy to check that for any $\alpha,\beta \in \F_{q^{n/2}}$, the element \[\alpha \gamma \zeta +\alpha \zeta x^{s(n-k)}+RF(x^n)\] is orthogonal to any element of $D$.
As a consequence, the set \[
\begin{split}
&\left\{a_0'\gamma \zeta+\sum\limits_{i=1}^{s(n-k)-1}a_ix^i+a_0''\zeta x^{s(n-k)}+RF(x^n) \colon a_i \in \F_{q^n},a_0',a_0'' \in \F_{q^{n/2}} \right\} \\
&= \zeta \gamma \left\{a_0'+\sum\limits_{i=1}^{s(n-k)-1}a_ix^i+a_0''\frac{1}{\gamma} x^{s(n-k)} + RF(x^n) \colon a_i \in \F_{q^n},a_0',a_0'' \in \F_{q^{n/2}} \right\} \\
&= \zeta \gamma D_{n,s,n-k}(1/\gamma,F)
\end{split}
\]
is contained in $D^{\perp}$ and by a dimensional argument we have that it coincides with $D^{\perp}$.

Now, let  \(
\C=\{M_{R_F}(a) \colon a \in D\} \subseteq M_n(\F_{q^s}). \)
The former computation shows that $\C^{\perp'}$ is equivalent to $D_{n,s,n-k}(1/\gamma,F)$ in $R_F \cong M_n(\F_{q^s})$. Finally, $\C^{\perp}=(\C^{\perp'})^{\top}$ is equivalent to to the adjoint of $D_{n,s,n-k}(1/\gamma,F)$ that is
\[
D_{n,s,n-k}(\sigma^{sk}(\gamma),\hat{F}) \subseteq R_{\hat{F}} \cong M_n(\F_{q^s})
\]
by \emph{\ref{adjointD})} of \Cref{prop:determinantionadjoint}, that proves our assertion.
\end{proof}

\end{document}
